The result reported here closes, at inter-city scale, the class of
frequency-ratio measurements that the SI-second redefinition process
depends on. Meeting the roadmap's frequency-ratio requirement has so
far required clocks on a shared campus, where a single stabilised
link serves both systems; here the two clocks sit in different cities,
in different institutions, and use different atomic species, and the
measurement still reaches the threshold, because the fibre link and
the two combs contribute well below the clocks themselves. The
bottleneck for reaching the redefinition threshold is therefore no
longer the transfer infrastructure but the systematic evaluations of
the clocks and the geodetic determination of their potential
difference.

The campaign also demonstrates the operational character that a
future optical-clock network will need. Over the two-month campaign
the two clocks were restarted, the transfer chain re-locked, and the
beat record interrupted many times; the segment-level ratio estimates
show no evidence of a systematic offset between the three measuring
stages, so the reproducibility of the comparison does not depend on
continuous operation of either system. This is the practical
counterpart of the coordinated multi-clock
campaigns~\citep{lindvall2025coordinated, pizzocaro2026european}: the
network tolerates interruption, provided each retained window is
self-consistently characterised.

Several limitations bound the present result. First, the per-segment
statistical uncertainties from the Allan-deviation extrapolation are
the dominant budget term, and five of the sixteen segments carry
model-limited fits; the random-effect treatment accounts for their
excess scatter at the level of the reported uncertainty, but a longer
record would reduce the statistical term. Second, the systematic
evaluations of the two clocks are those published by the
laboratories; independent verification of those budgets, as opposed
to the transfer, is not attempted here. Third, the geopotential
difference is taken from a first-order levelling campaign; an
alternative value exists in the experiment records and is carried as
an external item pending reconciliation, as is the physical-model
provenance of the supplied tide series and the exact accounting of
the full sample totals. Finally, the present campaign meets the
uncertainty requirement of the roadmap but not the full criterion,
which additionally requires each ratio to be measured at least twice
by different institutes and repeated over an extended period.

Looking ahead, the result establishes this link as a prototype node
pair of a distributed optical-clock network: the same infrastructure
can be extended to further nodes, and the technique demonstrated here
-- clocks at different sites agreeing at the redefinition threshold --
is the operation that both the redefinition and post-redefinition
applications require. Realising those applications over larger
distances will mean combining fibre networks of the kind used
here~\citep{schioppo2022cavities, chen2026dissemination} with
intercontinental techniques such as
VLBI~\citep{pizzocaro2021vlbi}, and will put the same demands on
geodetic infrastructure that this comparison
faced~\citep{grotti2018geodesy, mcgrew2018geodesy}.
